\id{ҒТАМР 65.65.03}{}

\begin{header}
\swa{М.Ч.~Тултабаев, Н.~Акжанов, А.~Сәдуакас, А.С.~Камали}{МАҚСАРЫ МАЙЫ НЕГІЗІНДЕ ОРТАША ЖӘНЕ ЖОҒАРЫ КАЛОРИЯЛЫ МАЙОНЕЗДІҢ РЕЦЕПТУРАСЫН ӘЗІРЛЕУ}

М.Ч.~Тултабаев,
Н.~Акжанов,
А.~Сәдуакас,
А.С.~Камали
\end{header}

\begin{affil}
«Қазақ қайта өңдеу және тағам өнеркәсіптері ҒЗИ» ЖШС, Астана, Қазақстан

\corrauthor{Корреспондент-автор: shomanyli@mail.ru}
\end{affil}

Мақалада мақсары майы негізінде орташа және жоғары калориялы майонез
рецептураларын әзірлеу және оларды ғылыми тұрғыдан негіздеу бойынша
жүргізілген кешенді зерттеудің нәтижелері баяндалады. Зерттеудің
өзектілігі тағам өнеркәсібіндегі жаңа үрдістермен, нақты айтұанда,
жасанды қоспаларды қолданбай-ақ физиологиялық тұрғыдан құнды өсімдік
майлары негізінде майлы эмульсиялық өнімдердің тағамдық және биологиялық
құндылығын арттыру қажеттілігімен айқындалады. Негізгі май фазасы
ретінде салқын сығымдау және өнеркәсіптік рафинация әдістерімен алынған,
полиқанықпаған май қышқылдарына, әсіресе линол қышқылына бай мақсары
майы пайдаланылды.

Зерттеу барысында мақсары майының май қышқылдық құрамы мен негізгі
физика-химиялық көрсеткіштері талданып, оның жоғары сапасы, тотығуға
төзімділігі және эмульсиялық тұздықтар өндірісінде қолдануға
технологиялық жарамдылығы дәлелденді. Алынған нәтижелер негізінде май
және су фазаларының әртүрлі арақатынасына ие орташа және жоғары
калориялы майонез рецептуралары әзірленді. Майонез өндірудің
технологиялық сұлбасы әзірленіп, ол тұрақты, микродисперсті эмульсияның
қалыптасуын қамтамасыз ететіні көрсетілді.

Дайын өнімдердің физика-химиялық көрсеткіштері бойынша жүргізілген
бағалау май мөлшері\-нің өзгеруі эмульсия тұрақтылығына теріс әсер
етпейтінін анықтады. Барлық зерттелген үлгілерде эмульсия тұрақтылығының
жоғары деңгейде сақталуы ұсынылған рецептуралардың технологиялық
сенімділігін дәлелдейді. Сонымен қатар, микробиологиялық зерттеулер
нәтижесінде барлық майонез үлгілерінің санитарлық-гигиеналық талаптарға
толық сәйкес келетіні анықталды.

Жүргізілген зерттеулердің нәтижелері мақсары майын орташа және жоғары
калориялы майонез өндірісінде қолданудың ғылыми және практикалық
тұрғыдан негізді екенін, сондай-ақ әзірленген рецептураларды
өнеркәсіптік деңгейде енгізуге болатынын көрсетеді.

{\bfseries Түйін сөздер:} майонез, калория, мақсары, физика-химиялық
қасиеттер, микрбиологиялық қауіпсіздік.

\begin{header}
РАЗРАБОТКА РЕЦЕПТА МАЙОНЕЗА СРЕДНЕЙ И ВЫСОКОЙ КАЛОРИЙНОСТИ НА ОСНОВЕ САФЛОРОВОГО МАСЛА

М.Ч.~Тултабаев\envelope,
Н.~Акжанов,
А.~Сәдуакас,
А.С.~Камали
\end{header}

\begin{affil}
ТОО «Казахский НИИ перерабатывающей и пищевой промышленности», Астана, Казахстан,

e-mail: shomanyli@mail.ru
\end{affil}

В статье рассматривается результаты проведенного комплексного
исследования по разработке и научному обоснованию рецептов майонеза
средней и высокой калорийности на основе сафлорового масла. Актуальность
исследования определяется современными тенденциями в пищевой
промышленности, в частности необходимостью повышения пищевой и
биологической ценности жирных эмульсионных продуктов на основе
физиологически ценных растительных масел без применения искусственных
добавок. В качестве основной масляной фазы использовалось сафлоровое
масло, богатое полиненасыщенными жирными кислотами, особенно линолевой
кислотой, полученное методами холодного отжима и промышленной рафинации.

В ходе работы были проанализированы жирнокислотный состав и основные
физико-химические показатели сафлорового масла, доказаны его хорошее
качество, устойчивость к окислению и технологическая пригодность для
применения в производстве эмульсионных соусов. На основе полученных
результатов разработаны рецептуры майонеза средней и высокой
калорийности с разными соотношениями жировой и водной фаз. Разработана
технологическая схема производства майонеза, которая обеспечивает
образование устойчивой микродисперсной эмульсии.

Проведенная оценка по физико-химическим показателям готовых продуктов
показала, что изменение количества масла не оказывает отрицательного
влияния на стабильность эмульсии. Сохранение на высоком уровне
стабильности эмульсии во всех исследованных образцах доказывает
технологическую стабильность предложенных рецептур. Кроме того, в
результате микробиологических анализов установлено, что все образцы
майонеза полностью соответствуют санитарно-гигиеническим требованиям.

Результаты проведенных исследований показывают, что применение
сафлорового масла в производстве майонеза средней и высокой калорийности
является научно и практически обоснованным, а также созданные рецептуры
могут быть внедрены на производственном уровне.

{\bfseries Ключевые слова:} майонез, калорийность, сафлор,
физико-химические свойства, микробиологическая безопасность.

\begin{header}
DEVELOPMENT OF A RECIPE FOR MEDIUM AND HIGH CALORIE MAYONNAISE BASED ON SAFFLOWER OIL

M.Ch.~Tultabayev\envelope,
N.~Akzhanov,
A Saduakas,
A.S.~Kamali
\end{header}

\begin{affil}
«Kazakh research Institute of processing and food industry» LLP, Astana, Kazakhstan,

e-mail: shomanyli@mail.ru
\end{affil}

The article presents the results of a comprehensive study on the
development of medium and high-calorie mayonnaise recipes based on
safflower oil and their scientific justification. The relevance of the
study is determined by new trends in the food industry, the need to
increase the nutritional and biological value of fatty emulsion products
based on physiologically valuable vegetable oils without the use of
artificial additives. Safflower oil, obtained by cold pressing and
industrial refining methods, rich in polyunsaturated fatty acids,
especially linoleic acid, was used as the main oil phase.

In the course of the study, the fatty acid composition and basic
physico-chemical indicators of safflower oil were analyzed, and its high
quality, oxidation resistance and technological suitability for use in
the production of emulsion sauces were proved. Based on the results
obtained, recipes for medium and high-calorie mayonnaise with different
ratios of fat and water phases were developed. A technological scheme
for the production of mayonnaise has been developed, and it has been
shown that it provides the formation of a stable, microdisperse
emulsion.

The assessment carried out according to the physico-chemical indicators
of finished products revealed that changes in the amount of oil do not
negatively affect the stability of the emulsion. Maintaining a high
level of emulsion stability in all studied samples proves the
technological reliability of the proposed formulations. In addition, as
a result of microbiological studies, it was found that all mayonnaise
samples fully comply with sanitary and hygienic requirements.

The results of the studies carried out indicate that the use of
safflower oil in the production of medium and high-calorie mayonnaise is
scientifically and practically justified, as well as that the developed
recipes can be implemented at an industrial level.

{\bfseries Keywords:} mayonnaise, caloric content, safflower,
physico-chemical properties, microbiological safety.

\begin{multicols}{2}
{\bfseries Кіріспе.} Майонез жоғары концентрациялы эмульсия ретінде
қарастырылады, ол негізінен өсімдік майы, қышқылдандырғыш және жұмыртқа
ұнтағынан тұрады және әлемдік тағам өнеркәсібінде жетекші орындардың
бірін алады. De Menezes R. C. F. және әріптестері қазіргі таңда дәстүрлі
жұмыртқа өнімдерін өсімдік компоненттерімен алмастыру, сондай-ақ
рафинирленген майлардың орнына физиологиялық тұрғыдан құнды майларды
қолдану тенденциясын атап көрсетеді. Әдеби шолу барысында инновациялық
ингредиенттердің негізгі түрлері жүйелендірілген: рапс олеосомалары, чиа
тұқымының майы, өсімдік камедилері, аквафаба, ақуыз изоляттары, өсімдік
ұнтақтары. Сондай-ақ майонездің стандартты құрамы көрсетілген: шамамен
63 \% май, 20 \% су, 8 \% тұрақтандырғыштар мен эмульгаторлар. Авторлар
алмастырғыштардың реологиялық қасиеттері рецептура компоненттерінің
қатынасына және өнімнің функционалды сипаттамаларына айтарлықтай әсер
ететінін атап өткен {[}1{]}.

Gao X. Және оның әріптестері өсімдік майының бір бөлігін тауық және
үйрек жұмыртқасының сарысына алмастыру эмульсияны біртекті әрі кішкене
бөлшекті етіп қалыптастыратынын, тұтқырлықты арттыратынын және
майонездің түсін жақсартатынын көрсетті. Сонымен қоса,
газ-хроматография, масс-спектрометрия әдістері үлгілердің дәмдік
профильдері арасындағы айырмашылықтарды растады {[}2{]}.

Технологиялық тұрғыдан тоқталсақ, Старовойтова К. В. және Терещук Л. В.
майонез өндірісінде әртүрлі пішінді жұмыртқа өнімдерін қолданудың
мән-жайын егжей-тегжейлі зерттеген. Олар фосфолипаза арқылы өңделген
құрғақ сарының ең аз мөлшері жұмыртқаның жалпы құрамын азайтуға
мүмкіндік беретінін, сонымен қатар өнімнің консистенциясы, эмульсияның
тұрақтылығы және крем тәрізді құрылымы бұзылмайтынын көрсетті {[}3{]}.
Karshenas M., Goli M. және Zamindar N. жұмыртқа сарысын арахис пен
кунжут шротынан алынған сүтпен бөліп айырбастауды ұсынды. Зерттеу
нәтижелері 50 \% алмастырылған жағдайда эмульсияның микроқұрылымы мен
май тамшыларының біртекті болатынын, ал 75--100 \% алмастырылғанда
реологиялық қасиеттердің нашарлайтынын көрсетті. Бұл зерттеулер өнімнің
функционалды құндылығын арттыру және экономикалық тиімділікті қамтамасыз
ету үшін басқа да ақуыз компоненттерін қолданудың мүмкіндігін дәлелдейді
{[}4{]}.

Майлы фазаның қалыптасуына және майонезді биологиялық активті заттармен
байытуға ерекше назар аударылған. Омельяненко О. С. және әріптестері
күнбағыс және өңделмеген қыша майларының 3:1 қатынасындағы рецептурасын
әзірлеп, «Провансаль» майонезінің дәстүрлі органолептикалық қасиеттерін
қалпына келтірді, май фазасының тотығуға төзімділігін арттырды және
өнімді витаминдер мен фитонцидтермен байытты {[}5{]}. Дурягина А. В.,
сондай-ақ Горнич Е. А. және әріптестері майонезге сүт альбуминын
енгізудің өнімнің биологиялық құндылығын арттыратынын, сонымен қатар
органолептикалық және физико-химиялық көрсеткіштерге әсер етпейтінін
көрсетті, бұл функционалды қоспалары бар жаңа рецептураларды әзірлеудің
маңыздылығын дәлелдейді. Резниченко И. Ю. және әріптестері күнбағыс
майының 10 \% мөлшерінде амарант майымен ішінара алмастырылуы тұтқырлық
пен эмульсия тұрақтылығын сақтауға, органолептикалық қасиеттерді
оңтауландыруға және скваленнің мөлшері арқасында биологиялық құндылықты
арттыруға мүмкіндік беретінін анықтады {[}6-8{]}.

Майлы фазаны функционализациялауда мақсары майы маңызды рөл атқарады.
Матеев Е. З. және әріптестері салқын сығымдау әдісімен алынған мақсары
майының құрамында омега-6 май қышқылдарының (шамамен 80 \% линолдық және
γ- линолдық) басым екенін және қышқылдық, пероксидтік және анизидиндік
көрсеткіштердің төмен екенін көрсетті, бұл оның жоғары сапасын және
майонез өндірісінде қолдануда болашағы зор екенін дәлелдейді {[}9{]}.
Teryokhina A. және әріптестері функционалды майонездерде май қышқылы
профилін мақсатты түрде модельдеуге, эмульсия тұрақтылығын бұзбай жүзеге
асыруға болатынын көрсетті {[}10{]}. Bredikhin S. A. және әріптестері
май фазасының құрамының және гомогенизация режимдерінің өнімнің
тұтқырлығы мен ағып кетуіне біршама әсер ететінін көрсетті {[}11{]}, ал
Gülen S. және әріптестері мақсары майын енгізу антирадикалды
белсенділікті, жалпы фенол құрамын және тұтынушыға тартымдылықты
арттыратынын, эмульсия тұрақтылығы мен микробиологиялық қауіпсіздіктің 6
ай бойы сақталатынын көрсетті {[}12{]}. Бұл мәліметтер орта мөлшердегі
калориялы майонезді жақсартылған тұрақтылықпен және ұзартылған сақтау
мерзімімен әзірлеуде мақсары майын қолданудың тиімділігін растайды.

{\bfseries Материалдар мен әдістер.} Зерттеу нысаны ретінде Қазақстан
Республикасының Жамбыл облысында өсірілген мақсары (Carthamus tinctorius
L.) майы өсімдігінің тұқымдары пайдаланылды. Мақсары майы «ҚазҚӨТӨ ҒЗИ»
ЖШС функционалдық тағам және тағам концентраттары зертханасында бірінші
суық сығымдау әдісімен алынды. Сонымен қатар, әртүрлі калориялық
деңгейдегі майонез рецептураларын әзірлеу үшін өнеркәсіптік жолмен
алынған рафинирленген, дезодорланған мақсары майы қолданылды.

Мақсары майы полиқанықпаған май қышқылдарының, әсіресе линол қышқылының
(70 \%-дан жоғары) жоғары мөлшерімен, сондай-ақ қышқылдық және
пероксидтік мөлшерінің төмен мәндерімен сипатталады, бұл оның
эмульсиялық тұздықтар, соның ішінде майонез өндірісінде қолдануға жоғары
биологиялық және технологиялық жарамдылығын көрсетеді.

Қосымша ингредиенттер ретінде майонез рецептурасына құрғақ сүт, қыша
ұнтағы, қант, ас тұзы, лимон шырыны, 80 \%-дық сірке қышқылы және ауыз
су енгізілді. Бұл компоненттердің нақты мөлшері әр рецептурада әртүрлі
өзгеріп, өнімнің дәмдік, қышқылдық және құрылымдық сипаттамаларын
реттеуге мүмкіндік береді

Физика-химиялық көрсеткіштерді анықтау келесі нормативтік құжаттарға
сәйкес жүргізілді:

- МЕМСТ 31762-2012 «Тамақ өнімдері. Ақуыздың, майдың, ылғалдың,
қышқылдық және пероксидтік санының массалық үлесін анықтау әдістері»;

- МЕМСТ ISO 3961-2020 «Жануар және өсімдік текті майлар мен тоңмайлар.
Йод санын анықтау»;

- МЕМСТ 30418-96 «Өсімдік және жануар майлары. Май қышқылдық құрамын
анықтау әдісі»;

- МЕМСТ 25555.0-82 «Тамақ өнімдері. Титрленетін қышқылдықты анықтау
әдісі».

{\bfseries Нәтижелер және талқылау.} Мақсары майы - құрамында
полиқанықпаған май қышқылдарына бай бағалы өсімдік майы. Оның
кұрылымында негізгі май қышқылдары ретінде линол қышқылы (C18:2), олеин
қышқылы (C18:1) және пальмитин қышқылы (C16:0) бар. Бұл май ағзаға
қажетті алмастырылмайтын май қышқылдарының зор көзі болып табылады,
жүрек-қан тамырлары жүйесінің қалыпты жұмысын қалыптастырады және
холестерин деңгейін төмендетуге оң әсер береді. Сонымен қатар, мақсары
майы тағамдық өндірістерде, соның ішінде майонез, салат сияқты дайын
өнімдерді дайындауда кеңінен пайдаланылады, себебі оның дәмдік және
құрылымдық қасиеттері жоғары (сур.1).

\fig[\columnwidth]{p2/image15}[1-сурет. Мақсары майының май қышқылдарының құрамы]

Диаграммада көрініп тұрғандай, мақсары майындағы май қышқылдарының ұлесі
айқын бөлінген. Линол қышқылы (C18:2) басымдыққа ие болып, май құрамының
шамамен 70\%-ын құрайды, бұл оның полиқанықпаған қасиетке бай екенін
және адам ағзасы үшін аса қажетті май қышқылы екенін көрсетеді. Олеин
қышқылы (C18:1) шамамен 20\% мөлшерде, бұл майға жақсы дәмдік қасиет
береді және жүрек-қан тамырлары жүйесінің қызметіне оң әсерін тигізеді.
Пальмитин қышқылы (C16:0) мен стеарин қышқылы (C18:0) кішкене мөлшерде -
тиісінше 6-7\% және 2-3\% - кездеседі, олардың болуы майдың құрылымын
тұрақтандыруға, эмульсиялық қасиеттерін жақсартуға ықпал етеді. Қалған
қышқылдардың үлесі өте аз (<1\%), бірақ олар да май сапасын
арттырып және оның функционалдық қасиеттерін сақтауға көмектеседі.

Мақсары майының физико-химиялық сипаттамаларын зерттеу оның сапасына
баға беріп, сақтау мерзімін анықтауға және тағамдық өнімдерде қолдану
жолдарын айқындауға мүмкіндік береді.1-кестеде майдың негізгі
физико-химиялық көрсеткіштері көрсетілген.
\end{multicols}

\tcap{1-кесте. Мақсары майының физико-химиялық көрсеткіштері}
\begin{longtblr}[
  label = none,
  entry = none,
]{
  cells = {c},
  cells = {font = \small},
  row{1} = {font=\bfseries\small},
  hlines,
  vlines,
}
Сипаттамалары & Өзгеру диапазоны & Эксперименталды мәндер\\
Салыстырмалы тығыздық & 0,919 - 0,924 & 0,920\\
Шағылу көрсеткіші, 25°С & 1,473 -1,476 & 1,474\\
Йод саны (г/100 г) & 115 - 155 & 125\\
Қышқылдық саны (мг КОН/г) & 0,78 - 5,76 & 1,87\\
Эфирлік саны (мг КОН/г) & - & 194,6\\
Пероксидтік саны & ≤10 & 9,3
\end{longtblr}

\begin{multicols}{2}
Салыстырмалы тығыздық 0,919-0,924 аралығында болып, эксперименталды мәні
0,920 құрайды, бұл майдың стандартты өсімдік майларының тығыздық
аралығына сәйкес екенін көрсетеді. Шағылу көрсеткіші 1,473-1,476
аралығында болып, 25°С-де 1,474 деп анықталды, бұл майдың жарықты сіңіру
және оптикалық қасиеттерін сипаттайды. Йод саны 115-155 г/100 г аралықта
болып, тәжірибелік мәні 125 г/100 г-ға тең, бұл майдың полиқанықпаған
май қышқылдарының көп екенін көрсетеді. Қышқылдық саны 0,78--5,76 мг
КОН/г аралығында өзгеріп, практикалық мәні 1,87 мг КОН/г болған, бұл
майдың сапасының жоғары екенін және қышқылдық деңгейінің төмендігін
білдіреді. Эфирлік саны 194,6 мг КОН/г құрап, май құрамындағы эфирлік
қосылыстарды бағалауға мүмкіндік береді. Пероксидтік саны 10-ға дейін
болып, тәжірибелік мәні 9,3 болған, бұл майдың тотығу процесіне
төзімділігін көрсетеді.

Жоғарыда сипатталған физико-химиялық қасиеттер негізінде майдың сапасы
мен қоректік құндылығы анықталғаннан кейін майонез рецептураларының
салыстырмалы құрамын көрсету маңызды.2-кестеде бақылау ``3 желания''
(күнбағыс майы), орташа калориялы және жоғары калориялы (мақсары майы)
майонездердің ингредиенттік үлестері көрсетілген. Бұл мәліметтер әр
рецептураның калориялылығын, май құрамын және су-құрғақ компонент
қатынасын бағалауға мүмкіндік береді.

Кестеден көріп отырғанымыздай, бақыланатын майонезінде май мөлшері ең
жоғары (65,4\%), су мөлшері салыстырмалы түрде төмен (24,25\%), бұл
классикалық, қою құрылымды майонезді береді. Орташа калориялы майонез
рецептурасында май мөлшері азайған (50,4\%), ал су мөлшері ұлғайған
(39,75\%), бұл майонезді жеңіл және калориясы орташа етеді. Жоғары
калориялы майонезде май мөлшері жоғарылап (65,4\%), су мөлшері аз
(23,75\%), нәтижесінде дәмі күшті, қою құрылымды өнім алынған.

Осы сипаттамаларды ескере отырып, мақсары майы негізіндегі майонез
үлгілерін алуға арналған технологиялық процесс әзірленді (сур.2).
\end{multicols}

\tcap{2-кесте. Жоғары және орташа калориялы майонез рецептуралары}
\begin{longtblr}[
  label = none,
  entry = none,
]{
  width = \linewidth,
  colspec = {Q[138]Q[312]Q[244]Q[246]},
  cells = {c},
  cells = {font = \small},
  row{1} = {font=\bfseries\small},
  hlines,
  vlines,
}
Ингредиент & Бақылау “3 желания майонезі” (күнбағыс майы) \% & Орташа калориялы (мақсары майы) \% & Жоғары калориялы (мақсары майы) \%\\
Май & 65,4 & 50,4 & 65,4\\
Жұмыртқа ұнтағы & 5,0 & 5,0 & 5,0\\
Сүт ұнтағы & 1,6 & 1,1 & 2,1\\
Қыша ұнтағы & 1,05 & 1,05 & 1,05\\
Қант & 1,05 & 1,05 & 1,05\\
Тұз & 1,1 & 1,1 & 1,1\\
80\% сірке қышқылы & 0,55 & 0,55 & 0,55\\
Су & 24,25 & 39,75 & 23,75
\end{longtblr}

\fig{p2/image16}[2-сурет. Мақсары майы негізіндегі майонез өндіру технологиясының кезеңдері]

\begin{multicols}{2}
Ұсынып отырған технологиялық сызбада шикізатты дайындау, сұйық фазаны
құрастыру, эмульсияны қалыптастыру, гомогенизация және pH түзету, дәмдік
қасиеттерді реттеу, сондай-ақ дайын өнімді орау және сақтау сатылары
жүйелі түрде бейнеленген. Бұл процестердің біркелкі орындалуы
эмульсияның тұрақтылығын арттырып, майонездің біртекті құрылымын,
қауіпсіздігін және жоғары сапалық көрсеткіштерін қамтамасыз етеді. Осы
технология негізінде өндірістік жағдайда қолдануға және функционалдық
өсімдік майлары негізіндегі жаңа өнімдер алуға бейімделген.

Дайын өнімдердің сапасын кешенді бағалау мақсатында бақыланатын және
тәжірибелік майонез үлгілерінің негізгі физика-химиялық көрсеткіштері
анықталды. Бұл көрсеткіштер рецептура құрамының эмульсия тұрақтылығына,
ылғалдылыққа және қышқылдық деңгейіне әсерін ғылыми тұрғыдан сипаттайды.
\end{multicols}

\tcap{3-кесте. Мақсары майы негізіндегі майонез үлгілерінің физика-химиялық көрсеткіштері}
\begin{longtblr}[
  label = none,
  entry = none,
]{
  width = \linewidth,
  colspec = {Q[267]Q[233]Q[219]Q[221]},
  cells = {c},
  cells = {font = \small},
  cell{1}{2} = {font=\bfseries\small},
  cell{1}{3} = {font=\bfseries\small},
  cell{1}{4} = {font=\bfseries\small},
  hlines,
  vlines,
}
\textbf{Көрсеткіш} & Бақылау “3 желания” (күнбағыс майы) & Орташа калориялы (мақсары майы) & Жоғары калориялы (мақсары майы)\\
Май мөлшері, \% & 65,4 & 50,4 & 65,4\\
Ылғал және ұшқыш заттар массалық үлесі, \% & 24,3 & 39,8 & 23,8\\
Қышқылдылығы (сірке қышқылына шаққанда), \% & 0,7 & 0,8 & 0,6\\
Эмульсия тұрақтылығы & 97 & 99 & 98
\end{longtblr}

\tcap{4-кесте. Дайын майонез үлгілерінің микробиологиялық көрсеткіштері}
\begin{longtblr}[
  label = none,
  entry = none,
]{
  width = \linewidth,
  colspec = {Q[229]Q[188]Q[150]Q[235]Q[137]},
  cells = {c},
  cells = {font = \small},
  row{1} = {font=\bfseries\small},
  hlines,
  vlines,
}
Көрсеткіш атауы & Көрсеткіштің реттелетін мәні & Бақылау “3 желания” & Орташа калориялы (мақсары майы) & Жоғары калориялы\\
МАжФАнМс, КТБ/см 3 (г), көп емес & ≤ 5,0×10³ & 95,41 & 54,6 & 60,2\\
ІТТБ & 0,1 г рұқсат етілмейді & анықталған жоқ & анықталған жоқ & анықталған жоқ\\
Ашытқылар мен саңырауқұлақ, КТБ/г & 50 г көп емес & анықталған жоқ & анықталған жоқ & анықталған жоқ\\
Salmonella spp. & 25 г аспауы тиіс & анықталған жоқ & анықталған жоқ & анықталған жоқ\\
Escherichia coli & 1г аспауы тиіс & анықталған жоқ & анықталған жоқ & анықталған жоқ\\
Staphylococcus aureus & 1 г аспауы тиіс & анықталған жоқ & анықталған жоқ & анықталған жоқ
\end{longtblr}

\begin{multicols}{2}
Кесте деректерінен сүйенсек, орташа калориялы үлгіде май мөлшерінің
төмендеуі ылғал үлесінің артуымен қатар жүреді, ал қышқылдылық пен
эмульсия тұрақтылығы нормативтік шектерде сақталған. Барлық үлгілерде
эмульсия тұрақтылығының 97--99 \% деңгейінде болуы майонез жүйелерінің
жоғары дисперстік тұрақтылығын және ұсынылған рецептуралардың
технологиялық тұрғыдан тиімді екенін айқындайды.

Дайындалған өнімдердің микробиологиялық қауіпсіздігі 4-кестеде
сипатталған.

Кестеден қарасақ, барлық дайын майонез үлгілері - бақылау (``3
желания'') және тәжірибелік (орташа және жоғары калориялы мақсары майы
негізіндегі) - микробиологиялық қауіпсіздік талаптарына толық сәйкес
келеді. МАжФАнМс көрсеткіші барлық үлгілерде ≤5,0×10³ КОЕ/г нормативінен
төменгі деңгейде болып, әсіресе тәжірибелік үлгілерде (54,6-60,2 КОЕ/г)
микробиологиялық тұрақтылықтың жоғары деңгейін көрсетеді, бұл майонездің
өндірістік технологиясының сапалы екенін дәлелдейді. ІТТБ (БГКП),
ашытқылар мен саңырауқұлақтар, сондай-ақ патогендік микроорганизмдер --
Salmonella spp., Escherichia coli және Staphylococcus aureus барлық
үлгілерде анықталмаған, яғни өнімдер санитарлық-гигиеналық нормаларға
толық сай келеді. Бұл сипаттамалар мақсары майы негізіндегі жаңа
рецептуралардың микробиологиялық тұрғыдан қауіпсіз және ұзақ мерзімге
сақталуға жарамды екенін, сондай-ақ олардың санитарлық-гигиеналық
стандарттармен үйлесетінін растайды. Сонымен қатар, тәжірибелік
үлгілердегі төмен МАжФАнМс деңгейі өнімнің органолептикалық және
физико-химиялық қасиеттерін сақтай отырып, ұзақ мерзімде тұрақтылығын
қамтамасыз ететінін көрсетеді.

Зерттеу нәтижелері бойынша орташа калориялы және жоғары калориялы
майонез үлгілері дайындалды. Орташа калориялы үлгідегі майонезде май
мөлшері азайғандықтан су үлесі ұлғайған, ал жоғары калориялы үлгілерде
май мөлшері жоғары болып, қою құрылымды өнім алынған. Микробиологиялық
зерттеулер көрсеткендей, барлық үлгілер санитарлық-гигиеналық талаптарға
сай келеді, эмульсия тұрақтылығы жоғары және физика-химиялық
көрсеткіштері нормативтік шектерде сақталған.

{\bfseries Қорытынды.} Мақсары майынан алынған орташа калориялы және жоғары
калориялы майонез үлгілері маңызды полиқанықпаған май қышқылдарының
мөлшерін көбейтеді, бұл өнімнің биологиялық құндылығын арттырады.
Физика-химиялық және микробиологиялық көрсеткіштердің нормативтік
стандартқа сәйкестігі, сондай-ақ эмульсиялардың тұрақтылығы мен
біркелкілігі эксперименттік құрамдардың қауіпсіз, технологиялық тиімді
және ұзақ мерзімді сақтауға жарамды екенін көрсетеді. Сондықтан мақсары
майына негізделген ұсынылған рецепт өнеркәсіптік жағдайларда сенімді
және функционалды майонез өндірісінде қолдану үшін болашағы мол екенін
көрсетеді.

{\bfseries \emph{Қаржыландыру}.} \emph{Жұмыс Қазақстан Республикасының Ауыл
шаруашылығы министрлігі BR22886613 "Мақсары майы негізінде функционалдық
мақсаттағы тағам өнімдерінің жаңа түрлерінің технологияларын әзірлеу"
тақырыбы бойынша "Ауыл шаруашылығы өсімдік шаруашылығының өнімі мен
шикізатын қайта өңдеу және сақтау жөніндегі инновациялық технологияларды
әзірлеу" қаржыландыратын бағдарлама шеңберінде жүргізілді.}
\end{multicols}

\begin{center}
{\bfseries Әдебиеттер}
\end{center}

\begin{refs}
1. De Menezes R.C.F., et al. Plant-based mayonnaise: Trending ingredients
for innovative products //International Journal of Gastronomy and Food
Science. -2022. -Vol.30(9): 100599. DOI
\href{https://doi.org/10.1016/j.ijgfs.2022.100599}{10.1016/j.ijgfs.2022.100599}.

2. Gao X., et al. Effect of adding different egg yolk oil on the flavor
and physicochemical properties of mayonnaise//Food Chemistry. -2025.
-Vol.477:143612. DOI
\href{https://doi.org/10.1016/j.foodchem.2025.143612}{10.1016/j.foodchem.2025.143612}.

3. Старовойтова К.В., Терещук Л.В. Разработка рецептур майонеза с учетом
основных тенденций совершенствования ассортимента //Техника и
технология пищевых производств. -2018. - Т.48(1). -С.91-98. DOI
10.21603/2074-9414-2018-1-91-98.

4. Karshenas M., Goli M., Zamindar N. Substitution of sesame and peanut
defatted-meal milk with egg yolk and evaluation of the rheological and
microstructural properties of low-cholesterol mayonnaise //Food
Science and Technology International. -2019. -Vol.25(8). -P.633-641.
DOI
~\href{https://doi.org/10.1177/1082013219853931}{10.1177/1082013219853931}.

5. Омельяненко О.С. и др. Разработка жировой основы майонеза, состоящей
из смеси растительных масел //Сборник трудов Российской национальной
научно-практической интернет-конференции для обучающихся и молодых
ученых 18-19 декабря 2019 г, Нижний Новгород -2020. -С.354-358.

6. Дурягина А. В. Разработка технологии майонеза с альбумином //Молодежь.
Наука. Инновации. -2021. -С.230-233.

7. Резниченко И. Ю., Захватаева А. В., Рубан Н. Ю. Разработка и оценка
качества майонезного соуса с использованием масла амаранта //Научный
журнал НИУ ИТМО. Серия «Процессы и аппараты пищевых производств».
-2022. -№2 (52). -С.3-9. DOI 10.17586-2310-1164-2022-15-2-3-9.

8. Горнич Е. А. и др. Разработка технологии майонеза, обогащённого
молочным сывороточным белком //Вестник АПК Верхневолжья. -2021.
-№1(53). -С.64-70. DOI 10.35694/YARCX.2021.53.1.011.

9. Матеев Е. З. и др. Исследование состава сафлорового масла, полученного
методом холодного прессования //Вестник Алматинского технологического
университета. -2019. -№ 2. -С.15-24.

10. Teryokhina A., Zheltoukhova E., Kopylov M. Analysis of fatty acid
composition of functional mayonnaise sauce with pumpkin products //BIO
Web of Conferences. -2024. -Vol.103:00075. DOI
\href{https://doi.org/10.1051/bioconf/202410300075}{10.1051/bioconf/202410300075}.

11. Bredikhin S.A., Martekha A.N., Andreev V.N. Research of the
rheological properties of mayonnaise with adding pumpkin and rice oils
to replace sunflower oil //Food Science and Technology. -2023. -Vol.
43: e67722. DOI https://doi.org/10.1590/fst.67722.

12. Gülen S., Günal-Köroğlu D., Turan S. Effects of cold-pressed oil
additives in varying proportions: Physico-chemical characteristics of
mayonnaises //Food Chemistry. -2025. --Vol.469: 142576.
DOI~\href{https://doi.org/10.1016/j.foodchem.2024.142576}{10.1016/j.foodchem.2024.142576}.
\end{refs}

\begin{center}
{\bfseries References}
\end{center}

\begin{refs}
1. De Menezes R.C.F., et al. Plant-based mayonnaise: Trending
ingredients for innovative products //International Journal of
Gastronomy and Food Science. -2022. -Vol.30(9): 100599. DOI
\href{https://doi.org/10.1016/j.ijgfs.2022.100599}{10.1016/j.ijgfs.2022.100599}.

2. Gao X., et al. Effect of adding different egg yolk oil on the flavor
and physicochemical properties of mayonnaise//Food Chemistry. -2025.
-Vol.477:143612. DOI
\href{https://doi.org/10.1016/j.foodchem.2025.143612}{10.1016/j.foodchem.2025.143612}.

3. Starovojtova K.V., Tereshhuk L.V. Razrabotka receptur majoneza s
uchetom osnovnyh tendencij sovershenstvovanija assortimenta //Tehnika i
tehnologija pishhevyh proizvodstv. -2018. - T.48(1). -S.91-98. DOI
10.21603/2074-9414-2018-1-91-98. {[}in Russian{]}

4. Karshenas M., Goli M., Zamindar N. Substitution of sesame and peanut
defatted-meal milk with egg yolk and evaluation of the rheological and
microstructural properties of low-cholesterol mayonnaise //Food Science
and Technology International. -2019. -Vol.25(8). -P.633-641. DOI
~\href{https://doi.org/10.1177/1082013219853931}{10.1177/1082013219853931}.

5. Omel' janenko O.S. i dr. Razrabotka zhirovoj osnovy
majoneza, sostojashhej iz smesi rastitel' nyh masel
//Sbornik trudov Rossijskoj nacional' noj
nauchno-prakticheskoj internet-konferencii dlja obuchajushhihsja i
molodyh uchenyh 18-19 dekabrja 2019 g, Nizhnij Novgorod -2020. -S.
354-358. {[}in Russian{]}

6. Durjagina A. V. Razrabotka tehnologii majoneza s
al' buminom //Molodezh'. Nauka. Innovacii.
-2021. -S.230-233. {[}in Russian{]}

7. Reznichenko I. Ju., Zahvataeva A. V., Ruban N. Ju. Razrabotka i ocenka
kachestva majoneznogo sousa s ispol' zovaniem masla
amaranta //Nauchnyj zhurnal NIU ITMO. Serija «Processy i apparaty
pishhevyh proizvodstv». -2022. -№2 (52). -S.3-9. DOI
10.17586-2310-1164-2022-15-2-3-9. {[}in Russian{]}

8. Gornich E. A. i dr. Razrabotka tehnologii majoneza, obogashhjonnogo
molochnym syvorotochnym belkom //Vestnik APK
Verhnevolzh' ja. -2021. -№1(53). -S.64-70. DOI
10.35694/YARCX.2021.53.1.011. {[}in Russian{]}

9. Mateev E. Z. i dr. Issledovanie sostava saflorovogo masla,
poluchennogo metodom holodnogo pressovanija //Vestnik Almatinskogo
tehnologicheskogo universiteta. -2019. -№ 2. -S.15-24. {[}in Russian{]}

10. Teryokhina A., Zheltoukhova E., Kopylov M. Analysis of fatty acid
composition of functional mayonnaise sauce with pumpkin products //BIO
Web of Conferences. -2024. -Vol.103:00075. DOI
\href{https://doi.org/10.1051/bioconf/202410300075}{10.1051/bioconf/202410300075}.

11. Bredikhin S.A., Martekha A.N., Andreev V.N. Research of the
rheological properties of mayonnaise with adding pumpkin and rice oils
to replace sunflower oil //Food Science and Technology. -2023. -Vol.43:
e67722. DOI https://doi.org/10.1590/fst.67722.

12. Gülen S., Günal-Köroğlu D., Turan S. Effects of cold-pressed oil
additives in varying proportions: Physico-chemical characteristics of
mayonnaises //Food Chemistry. -2025. --Vol.469: 142576.
DOI~\href{https://doi.org/10.1016/j.foodchem.2024.142576}{10.1016/j.foodchem.2024.142576}.
\end{refs}

\begin{info}
\hspace{1em}\emph{{\bfseries Авторлар туралы мәліметтер}}

Тултабаев М.Ч - техника ғылымдарының докторы, «Қазақ қайта өңдеу және
тағам өнеркәсіптері ғылыми-зерттеу институты» ЖШС Астана филиалы,
Астана, Қазақстан, e-mail: shomanyli@mail.ru;

Акжанов Н. - PhD докторант, «Қазақ қайта өңдеу және тағам
өнеркәсіптері ғылыми-зерттеу институты» ЖШС Астана филиалы, Астана,
Қазақстан, e-mail: nurtore0308@gmail.com;

Садуакас А. - «Қазақ қайта өңдеу және тағам өнеркәсіптері
ғылыми-зерттеу институты» ЖШС Астана филиалы, Астана, Қазақстан,
e-mail: aykon96@mail.ru;

Камали А.С. - «Қазақ қайта өңдеу және тағам өнеркәсіптері
ғылыми-зерттеу институты» ЖШС Астана филиалы, Астана,Қазақстан,
e-mail: akerke.kamali@bk.ru;

\hspace{1em}\emph{{\bfseries Information about the authors}}

Tultabayev M.Ch.- Doctor of Technical Sciences, Astana branch of
Kazakh research Institute of processing and food industry LLP, Astana,
Kazakhstan, e-mail: shomanyli@mail.ru;

Akzhanov N.- PhD doctoral student, Astana branch of Kazakh research
Institute of processing and food industry LLP, Astana, Kazakhstan,
e-mail: nurtore0308@gmail.com;

Saduakas A.- Astana Branch of the Kazakh Research Institute of
Processing and Food Industries, Astana, Kazakhstan, e-mail:
aykon96@mail.ru;

Kamali A.S. - Astana Branch of the Kazakh Research Institute of
Processing and Food Industries, Astana, Kazakhstan, e-mail:
akerke.kamali@bk.ru.
\end{info}
